\documentclass{beamer}
\usetheme{Madrid}
\usepackage{amsmath, amssymb, amsthm}
\usepackage{graphicx}
\usepackage{listings}
\usepackage{gensymb}
\usepackage[utf8]{inputenc}
\usepackage{hyperref}
\usepackage{gvv}
\usepackage{tikz}
\lstset{
  language=Python,
  basicstyle=\ttfamily\small,
  keywordstyle=\color{blue},
  stringstyle=\color{orange},
  numbers=left,
  numberstyle=\tiny\color{gray},
  breaklines=true,
  showstringspaces=false
}
\usetikzlibrary{decorations.pathmorphing}

\title{Question-5.4.17}
\author{EE25BTECH11022 - sankeerthan}
\date{}
\begin{document}

\frame{\titlepage}

\begin{frame}
\frametitle{Question}
Using elementary transformations, find the inverse of the following matrix. 
\begin{align*}
\myvec{3&1\\5&2}
\end{align*}
\end{frame}

\begin{frame}
\frametitle{Solution}
Given  
\begin{align}
\vec{A}=\myvec{3&1\\5&2}
\end{align}
Let $\vec{A}^{-1}$ be the inverse of $\vec{A}$.Then
\begin{align}
    \vec{A}\vec{A}^{-1}=\vec{I}
\end{align}
Augmented matrix of $\augvec{1}{1}{\vec{A} & \vec{I}}$ is given by
\begin{align}
\augvec{2}{2}{3 & 1 & 1 & 0 \\ 5 & 2 & 0 & 1} \xrightarrow{R_2 \rightarrow 3R_2 - 5R_1}\augvec{2}{2}{3 & 1 & 1 & 0 \\ 0 & 1 & -5 & 3} \xrightarrow{R_1 \rightarrow R_1 - R_2}\augvec{2}{2}{3 & 0 & 6 & -3 \\ 0 & 1 & -5 & 3} 
\end{align}
\begin{align}
\augvec{2}{2}{3 & 0 & 6 & -3 \\ 0 & 1 & -5 & 3} \xrightarrow{R_1 \rightarrow \tfrac{1}{3} R_1}\augvec{2}{2}{1 & 0 & 2 & -1 \\ 0 & 1 & -5 & 3}
\end{align}
Hence the inverse of the matrix $\myvec{3&1\\5&2}$ is \myvec{2&-1\\-5&3}
\end{frame}

\begin{frame}[fragile]
\frametitle{C-Code}
\begin{lstlisting}[language=C]
// code.c
#include <stdio.h>

// Function to multiply matrix and vector
void mat_vec_mult(double mat[2][2], double vec[2], double result[2]) {
    for (int i = 0; i < 2; ++i) {
        result[i] = 0;
        for (int j = 0; j < 2; ++j) {
            result[i] += mat[i][j] * vec[j];
        }
    }
}
\end{lstlisting}
\end{frame}
\begin{frame}[fragile]
\frametitle{C-Code}
\begin{lstlisting}[language=C]
// Solves the equation X = A_inv * B where A_inv is the inverse of the given matrix
void find_points(double x_vals[2], double y_vals[2], double points[2][2]) {
    // Inverse of given matrix A [[3, 1], [5, 2]] is [[2, -1], [-5, 3]]
    double inv[2][2] = {{2, -1}, {-5, 3}};
    for (int i = 0; i < 2; ++i) {
        double b[2] = {x_vals[i], y_vals[i]};
        double res[2];
        mat_vec_mult(inv, b, res);
        points[i][0] = res[0];
        points[i][1] = res[1];
    }
}
\end{lstlisting}
\end{frame}
\begin{frame}[fragile]
\frametitle{C-Code}
\begin{lstlisting}[language=C]
// Main just prints out by default
int main() {
    double x_vals[2] = {1, 4}; // Example points on x
    double y_vals[2] = {1, 4}; // Example points on y
    double points[2][2];
    find_points(x_vals, y_vals, points);
    for (int i = 0; i < 2; ++i) {
        printf("%lf %lf\n", points[i][0], points[i][1]);
    }
    return 0;
}
\end{lstlisting}
\end{frame}
\begin{frame}[fragile]
\frametitle{Python-Code}
\begin{lstlisting}
# Code by GVV Sharma
# Modified for Problem Solution
# Released under GNU GPL
# Calculating area enclosed between curves
# call.py
import ctypes
import numpy as np
import os

so_file = os.path.abspath('./code.so')
lib = ctypes.CDLL(so_file)

\end{lstlisting}
\end{frame}
\begin{frame}[fragile]
\frametitle{Python-Code}
\begin{lstlisting}
# define C double arrays
class DoubleArray2(ctypes.Structure):
    _fields_ = [("arr", ctypes.c_double * 2)]
def create_double2(np_array):
    array_type = ctypes.c_double * 2
    return array_type(*np_array)

# C function
double2 = ctypes.c_double * 2
points_type = double2 * 2

find_points = lib.find_points
find_points.argtypes = [double2, double2, points_type]
find_points.restype = None
\end{lstlisting}
\end{frame}

\begin{frame}[fragile]
\frametitle{Python-Code}
\begin{lstlisting}
def get_points(x_vals, y_vals):
    x_c = double2(*x_vals)
    y_c = double2(*y_vals)
    out = points_type()
    find_points(x_c, y_c, out)
    result = np.array([[out[i][0], out[i][1]] for i in range(2)])
    return result

if __name__ == '__main__':
    x_vals = [1.0, 4.0]
    y_vals = [1.0, 4.0]
    pts = get_points(x_vals, y_vals)
    print(pts)
\end{lstlisting}
\end{frame}

\end{document}